\begin{example}[CNOT gate]\label{exa:cnotBenc}
    The CNOT gate is the controlled Pauli-X, that is
    \begin{align*}
        \mathrm{CNOT}\left[\cheadvariable{\insymbol},\cheadvariable{\outsymbol},\catvariable\right]
        = \paulixat{\cheadvariable{\insymbol},\cheadvariable{\outsymbol}} \otimes \tbasisat{\catvariable}
        + \identityat{\cheadvariable{\insymbol},\cheadvariable{\outsymbol}} \otimes \fbasisat{\catvariable} \, .
    \end{align*}
    For the identity function
    \begin{align*}
        \mathrm{Id}(x) = x \, .
    \end{align*}
    we have
    \begin{align*}
        \qcbencodingofat{\mathrm{Id}}{\cheadvariable{\insymbol},\cheadvariable{\outsymbol},\catvariable}
        = \mathrm{CNOT}\left[\cheadvariable{\insymbol},\cheadvariable{\outsymbol},\catvariable\right] \, .
    \end{align*}
    The $\mathrm{CNOT}$ is therefore the \computationCircuit{} to the identity function.
\end{example}